\begin{abstract} % unterdrueckt u.a. Kopf- und Fusszeile
\begin{center}

\begin{minipage}{12cm}

    {\Large\textit{Kurzzusammenfassung}}

    Die Problemstellung \bzw die Motivation dieser Arbeit ist es aufzuzeigen wie das Regelwerk des erfolgreichsten deutschen Stift-und-Papier-Rollenspiels \AbkN{DSA} in der aktuellen 5ten-Ausgabe so formalisiert werden kann, dass ein möglichst wartbares und quelloffenes Hilfsmittel für die Charakter-Erstellung und perspektivisch eine Basis für die Charakter-Verwaltung entsteht.
    Das Regelwerk zeichnet sich \insb durch den Umfang und die Komplexität der Regeln und Kombinationen der Charakter-Merkmale und Merkmalsausprägungen aus sowie der Vielzahl an Regelwerkserweiterungen.
    Aufgrund dessen hat die Spielenden-Gemeinschaft Hilfsmittel entwickelt, um die Aufgabe der Charakter-Erstellung und -Verwaltung zu vereinfachen.
    Keines der bisher verbreiteten Hilfsmittel, welches Charakter-Validierungen durchführt, unterstützt die aktuellste Ausgabe.
    Oft wurde als Grund die weitreichenden notwendigen Anpassungen an der Software und die fehlenden Kapazitäten dafür genannt.
    In dieser Arbeit werden zunächst die Rahmenbedingungen für quelloffene Projekte erörtert und ein geeignetes Formalisierungsverfahren samt Framework und Programmiersprache ermittelt.
    Anschließend wird ein \textquote{Kochbuch} aufgestellt, wie bestimmte Arten von Regeln des Regelwerks, unter festgelegten Konventionen, formalisiert werden könnten.
    Abschließend wird dieses Konzept als Proof-of-Concept umgesetzt und beides kritisch bewertet.
    Es wird zum Ergebnis kommen, dass eine Formalisierung möglich und einfach ist sowie die gut durchdachte Integration dessen in eine Software als Hilfsmittel eine wartbare quelloffene Lösung für die Gemeinschaft ist.

%Ersteres kann ein Indiz für eine zu starke Verflechtung der fachlichen Logik (das Regelwerk) mit der technischen Implementierung sein, sodass diese Softwares sehr speziell sind.
%Letzteres könnte mehrere Gründe haben, wie eine schlechte Wartbarkeit und große Komplexität oder eine nicht quelloffenen Software.
%    Diese Probleme werden in dieser Arbeit angeganen

\end{minipage}
\end{center}
\end{abstract}

\clearpage

%    \vspace{5em}
\begin{abstract} % unterdrueckt u.a. Kopf- und Fusszeile
\begin{center}
\begin{minipage}{12cm}

    {\Large\textit{Abstract}}

    The problem and motivation of this work is to demonstrate how the rulebook of the most popular pen-and-paper role-playing game in Germany, \AbkN{DSA}, can be formalized in its latest 5th edition, creating a maintainable, open-source tool for character creation and ultimately also character management.
    The rulebook is known for its complexity by the combination of character traits and trait characteristics as well as the extensive number of rulebook extensions.
    As a result, the community has developed tools to simplify the task of character creation and management.
    However, none of the widely used tools, doing character validation, support the latest edition due to self-reported extensive required customizations and lack of development capacity to tackle this.
    This paper discusses the general conditions for open-source projects and identifies a suitable formalization procedure, including the framework and programming language.
    It provides a \textquote{cookbook} outlining how to formalize different rules of the rulebook under specified conventions.
    Finally, the paper implements a proof-of-concept and evaluates both critically.
    It concludes that a formalization is achievable and quite simple, and the well-designed integration of this into a software created a great maintainable and open-source solution for the community.

\end{minipage}
\end{center}
\end{abstract}